\documentclass[12pt,a4paper]{article}

\usepackage[T1]{fontenc}
\usepackage[utf8]{inputenc}
\usepackage[romanian]{babel}
\usepackage{geometry}

\geometry{margin=2.5cm}

\title{Tutoriat 2\\Const și Static}
\author{}
\date{}

\begin{document}

\maketitle

\section{Keyword-ul const}

În proiectarea claselor există situații în care anumite date nu trebuie modificate după inițializare. Pentru exprimarea acestei proprietăți, C++ oferă calificatorul const.

Acesta reprezintă o promisiune făcută compilatorului că valoarea respectivă nu va fi modificată pe durata de viață a obiectului.

\subsection{Exemple de utilizare}

\begin{itemize}
\item Clasa Student: câmpul dataNasterii.
\item Clasa CardBancar: câmpul numarCard.
\end{itemize}

\subsection{Variabile constante}

Exemplu:

const int x = 3;

const double y = 3.0;

Orice încercare de modificare a unei variabile constante produce eroare de compilare.

\textbf{Observație}

O variabilă constantă trebuie inițializată deoarece ulterior nu mai poate primi o valoare prin operatorul de atribuire.

\section{Pointeri constanți vs pointeri către constante}

Exemplu:

const int* p;

Se citește de la dreapta la stânga:

pointer către un întreg constant.

Exemplu:

p = \&x;      // permis

*p = 5;      // interzis

Exemplu:

int* const p = \&x;

Se citește:

pointer constant către un întreg.

Exemplu:

p = \&y;      // interzis

*p = 5;      // permis

\section{Date membre constante}

Membrii constanți ai unei clase trebuie inițializați prin lista de inițializare a constructorului.

\subsection{Lista de inițializare}

Exemplu:

class Student \{

    const string CNP;

    int varsta;

public:

    Student(string \_CNP)
        : CNP(\_CNP)
    \{
        varsta = 0;
    \}
\};

Lista de inițializare este introdusă prin caracterul ":" și este executată înaintea corpului constructorului.

Pentru datele constante aceasta este obligatorie.

\section{Metode constante}

O metodă constantă garantează faptul că obiectul apelant nu va fi modificat.

Exemplu:

class Student \{

    int varsta;

public:

    int getVarsta() const \{
        return varsta;
    \}
\};

În interiorul unei metode constante nu este permisă modificarea datelor membre.

Exemplu:

void set() const \{

    varsta = 30; // eroare

\}

Getter-ele sunt aproape întotdeauna declarate const.

\section{Erori frecvente provocate de const}

\subsection{Obiecte constante}

Un obiect constant poate apela exclusiv metode constante.

Exemplu:

class Student \{

    int varsta;

public:

    int getVarsta() \{
        return varsta;
    \}
\};

int main() \{

    const Student s;

    cout << s.getVarsta(); // eroare
\}

Varianta corectă:

int getVarsta() const \{
    return varsta;
\}

\subsection{Referințe neconstante}

Exemplu:

void sum(Student\& s)

acceptă modificarea parametrului.

Pentru obiecte constante trebuie folosit:

void sum(const Student\& s)

Acest lucru garantează compilatorului că argumentul nu va fi modificat.

\section{Keyword-ul static}

\subsection{În limbajul C}

Exemplu:

void f() \{

    static int nr = 0;

    nr++;

    printf("\%d", nr);
\}

Variabila este inițializată o singură dată și își păstrează valoarea între apelurile funcției.

\section{Membri statici ai unei clase}

Membrii statici nu aparțin unui obiect, ci întregii clase.

Toate instanțele clasei partajează aceeași valoare.

\subsection{Date statice neconstante}

Exemplu:

class A \{

public:

    static int x;
\};

int A::x = 3;

Linia

int A::x = 3;

este obligatorie.

\subsection{Date statice constante}

Exemplu:

class A \{

public:

    const static int x = 3;
\};

Pentru datele statice constante inițializarea se face direct în clasă.

\section{Metode statice}

Exemplu:
\\
class A \{

public:

    static void m1() \{
        cout << "m1";
    \}

    static void m2();\\
\};
\\
void A::m2() \{
    cout << "m2";
\}

O metodă statică poate fi apelată prin:

A::m1();

sau

A a;

a.m2();

deși prima variantă este recomandată.

\textbf{Regulă importantă}

O metodă statică poate accesa exclusiv membri și metode statice.

\section{Erori frecvente provocate de static}

\subsection{Dată statică neinițializată}

Exemplu:
\\
class A \{

public:

    static int x;\\
\};

int main() \{

    cout << A::x; // eroare\\
\}

Trebuie adăugată definiția:

int A::x = 3;

\subsection{Accesarea membrilor nestatici}

Exemplu:
\\
class A \{

public:

    int x = 3;

    static void m() \{

        cout << x; // eroare\\
    \}\\
\};

Rezolvări:

\begin{itemize}
\item transformarea lui x în membru static;
\item eliminarea calificatorului static din metodă;
\item crearea unui obiect local și accesarea prin instanță.
\end{itemize}

\section{Concluzii}

\begin{itemize}
\item const descrie imutabilitatea și oferă garanții la compilare;
\item membrii constanți folosesc lista de inițializare;
\item obiectele constante pot apela doar metode constante;
\item membrii statici aparțin clasei, nu instanțelor;
\item metodele statice au acces exclusiv la membri statici.
\end{itemize}

\end{document}